% =====================================================================
% Primera parte – 2. Propuesta de solución
% =====================================================================
\section{Propuesta de solución}

\subsection{Objetivos}

\paragraph{Objetivo general.} Desarrollar una aplicación web que le permita a una persona registrar y clasificar sus ingresos y gastos, controlar la ejecución de sus presupuestos por categoría y período, y cumplir metas de ahorro mediante un plan diario con indicador de progreso.

\paragraph{Objetivos específicos.}
\begin{enumerate}[label=OE\arabic*., leftmargin=1.6cm]
  \item Reducir el registro de un movimiento a un formulario único con validaciones en el cliente y en el servidor (ataca C1).
  \item Calcular automáticamente el gasto ejecutado y el porcentaje real de cada presupuesto, y alertar cuando se supera el umbral configurado (ataca C2).
  \item Generar sin intervención del usuario las transacciones recurrentes del día (ataca C3).
  \item Convertir cada meta de ahorro en un plan de cuotas diarias sin valores negativos y mostrar su avance (ataca C4).
  \item Verificar la calidad del producto con pruebas unitarias, de integración, de sistema y de aceptación, y con métricas basadas en ISO/IEC 25010 \parencite{iso25010}.
\end{enumerate}

\subsection{Usuarios y actores}

\begin{table}[H]
\caption{Actores del sistema}
\label{tab:actores}
\begin{tabularx}{\linewidth}{L{3.6cm} L{2.2cm} Y}
\toprule
Actor & Tipo & Descripción \\
\midrule
Usuario & Principal, humano & Persona que registra sus finanzas: se registra o inicia sesión y usa todos los módulos. \\
Programador de tareas & Secundario, tiempo & Proceso de APScheduler que cada día a las 00:05 genera las transacciones recurrentes. \\
Google OAuth & Secundario, sistema externo & Valida la identidad del usuario cuando inicia sesión con su cuenta de Google. \\
\bottomrule
\end{tabularx}
\tablenote{Fuente: elaboración propia.}
\end{table}

\subsection{Descripción del sistema}\label{sec:descripcion}

ShiftCash es una aplicación web cliente-servidor con dos subsistemas desacoplados que se comunican por HTTPS con una API REST/JSON:

\begin{itemize}
  \item \textbf{Frontend} (\codigo{frontend/src}): aplicación de página única en React 19 y Vite \parencite{react2025,vite2025}, desplegada en Vercel. Usa React Router con las rutas públicas \codigo{/}, \codigo{/login}, \codigo{/registro} y la ruta privada \codigo{/dashboard}, protegida por \codigo{PrivateRoute}. Dentro del dashboard se navega entre siete vistas. Se organiza en capas: \codigo{api/} (clientes HTTP con Axios, que adjuntan el JWT), \codigo{hooks/} (estado y reglas de presentación), \codigo{components/} (vistas, modales y diseño) y \codigo{pages/}.
  \item \textbf{Backend} (\codigo{backend/app}): API REST en FastAPI \parencite{fastapi2025} desplegada en Render. Se organiza en capas: \codigo{models.py} (entidades SQLAlchemy), \codigo{schemas.py} (validación con Pydantic), \codigo{services/} (reglas de negocio), \codigo{routes/} (endpoints delgados que delegan en los servicios), \codigo{auth.py} (JWT y bcrypt) y \codigo{scheduler.py} (tarea diaria de recurrencias).
  \item \textbf{Base de datos}: PostgreSQL alojado en Supabase, al que la API accede con SQLAlchemy y \codigo{psycopg2} mediante la variable de entorno \codigo{DATABASE_URL}.
\end{itemize}

\subsubsection{Módulos y funciones}

\begin{longtable}{L{2.6cm} L{7.2cm} L{2.4cm}}
\caption{Módulos de ShiftCash y su relación con los modelos}\label{tab:modulos}\\
\toprule
Módulo & Funciones & RF / CU \\
\midrule
\endfirsthead
\toprule
Módulo & Funciones & RF / CU \\
\midrule
\endhead
Autenticación & Registro con correo y contraseña; inicio de sesión con contraseña o con Google; emisión y verificación del JWT; protección de rutas. & RF01–RF02 / CU04–CU06 \\
Transacciones & Registrar, editar y eliminar ingresos y gastos; marcar una transacción como recurrente; generación diaria automática de las recurrentes. & RF03–RF06 / CU01, CU08, CU14 \\
Historial & Listado agrupado por fecha con búsqueda y filtros por tipo, categoría, recurrencia, frecuencia y fechas; detalle de cada transacción. & RF07 / CU09 \\
Resumen y gráficas & Totales del mes (gastado, ganado, ahorrado, número de transacciones), alertas de presupuesto y gráficas de torta, barras y áreas \parencite{recharts2025}. & RF08, RF09, RF14 / CU10, CU11 \\
Categorías & Categorías predeterminadas (no editables) y personalizadas (crear, editar, eliminar si no tienen transacciones). & RF10 / CU12 \\
Presupuestos & Límite por categoría y período (diario, semanal, mensual o único, permanente o no); gasto ejecutado, porcentaje real y estado «superado». & RF11–RF13 / CU02 \\
Metas de ahorro & Meta con objetivo y fechas; plan de cuotas diarias redondeadas a la moneda; marcado de días y progreso. & RF15–RF18 / CU03, CU13 \\
Perfil y configuración & Nombre, correo, contraseña, moneda, foto, umbral de alerta, tema claro u oscuro y eliminación de la cuenta. & RF19, RF20 / CU07 \\
\bottomrule
\end{longtable}
\vspace{-1em}
{\small\textit{Nota.} Los RF se detallan en la Tabla \ref{tab:rf} y los CU en la sección \ref{sec:cu}. Fuente: elaboración propia.}

\subsubsection{Flujo principal}

El uso típico de ShiftCash sigue la secuencia de la Figura \ref{fig:flujo}: el usuario inicia sesión, define sus presupuestos y metas, registra sus movimientos (o deja que las recurrentes se generen solas) y consulta en el resumen cuánto lleva gastado y ahorrado. Cada paso corresponde a uno o más casos de uso de la sección \ref{sec:cu}.

\begin{figure}[H]
\centering
\caption{Flujo principal de uso de ShiftCash}
\label{fig:flujo}
\resizebox{\linewidth}{!}{%
\begin{tikzpicture}[node distance=5mm and 6mm, paso/.style={caja, text width=2.5cm, minimum height=1.4cm}]
\node[paso] (a) {Iniciar sesión\\(CU05/CU06)};
\node[paso, right=of a] (b) {Crear presupuestos y metas\\(CU02, CU03)};
\node[paso, right=of b, cajaacento] (c) {Registrar ingresos y gastos\\(CU01)};
\node[paso, right=of c] (d) {Calcular gasto, porcentaje y progreso};
\node[paso, right=of d] (e) {Consultar resumen, historial y gráficas\\(CU09–CU11)};
\node[paso, cajagris, below=8mm of c] (r) {Generar recurrentes 00:05\\(CU14)};
\draw[flecha] (a) -- (b); \draw[flecha] (b) -- (c); \draw[flecha] (c) -- (d); \draw[flecha] (d) -- (e);
\draw[flecha] (r) -- (c);
\draw[flechapunteada] (e.south) |- ++(0,-2.4) -| node[pos=0.25, below, font=\scriptsize]{ajustar presupuestos o metas} (b.south);
\end{tikzpicture}}
\figurenote{Fuente: elaboración propia.}
\end{figure}

\subsubsection{Procesos automatizados}

\begin{table}[H]
\caption{Procesos que ShiftCash automatiza}
\label{tab:procesos}
\begin{tabularx}{\linewidth}{L{3.4cm} Y L{3.6cm}}
\toprule
Proceso & Qué hace & Dónde está \\
\midrule
Transacciones recurrentes & Cada día a las 00:05 replica las recurrentes (mensual, semanal, quincenal, anual) que correspondan a la fecha, sin duplicarlas. & \codigo{scheduler.py} \\
Gasto de un presupuesto & Calcula el rango del período, suma los gastos de la categoría y devuelve el porcentaje real (puede superar 100\,\%). & \codigo{services/budget_service.py} \\
Plan de ahorro & Reparte el objetivo en cuotas diarias variables (±30\,\%) redondeadas a la unidad de la moneda del usuario; ninguna es negativa y la suma es exacta. & \codigo{services/saving_service.py} \\
Progreso de la meta & Total ahorrado, porcentaje y días cumplidos; se rehace el plan si cambian el objetivo o la fecha final. & \codigo{services/saving_service.py} \\
Resumen y alertas & Totales del mes y alerta cuando el gasto supera el umbral configurado (80\,\% por defecto). & \codigo{useFinanzas.js}, \codigo{VistaResumen.jsx} \\
Integración continua & En cada push: lint y compilación del frontend y pruebas del backend con cobertura. & \codigo{.github/workflows} \\
\bottomrule
\end{tabularx}
\tablenote{Fuente: elaboración propia.}
\end{table}

\subsubsection{Tecnologías}

\begin{table}[H]
\caption{Tecnologías utilizadas}
\label{tab:tecnologias}
\begin{tabularx}{\linewidth}{L{2.6cm} L{4.6cm} Y}
\toprule
Capa & Tecnología (versión) & Uso \\
\midrule
Interfaz & React 19.2, Vite 8.0, React Router 7.14 & SPA, empaquetado y rutas públicas y privadas. \\
Gráficas y HTTP & Recharts 3.8, Axios 1.16 & Gráficas del resumen; cliente HTTP con el JWT. \\
Identidad & \codigo{@react-oauth/google} 0.13 & Inicio de sesión con Google. \\
API & FastAPI 0.136, Uvicorn 0.49 & Endpoints REST y documentación OpenAPI. \\
Dominio y datos & SQLAlchemy 2.0, Pydantic 2.13 & ORM y validación de entradas \parencite{bayer2025,pydantic2025}. \\
Seguridad & python-jose 3.5, bcrypt 5.0 & JWT y hash de contraseñas. \\
Tareas & APScheduler 3.11 & Recurrencias diarias. \\
Base de datos & PostgreSQL (Supabase) & Persistencia. \\
Despliegue & Vercel, Render, GitHub Actions & Hospedaje e integración continua. \\
Pruebas & pytest 9.2, pytest-cov 7.1, Cypress, Selenium 4, JMeter & Pruebas y medición (Segunda y Tercera parte). \\
\bottomrule
\end{tabularx}
\tablenote{Versiones tomadas de \codigo{package.json} y \codigo{backend/requirements.txt}. Fuente: elaboración propia.}
\end{table}

\paragraph{Relación con los modelos.} Cada módulo se traza hasta los requisitos (Tabla \ref{tab:rf}), los casos de uso (Tabla \ref{tab:trazabilidad}) y las clases que lo implementan (sección \ref{sec:clases-impl}); los servicios del backend son los mismos objetos que reciben los mensajes en los diagramas de secuencia (sección \ref{sec:secuencia}).

\subsection{Alcance}\label{sec:alcance}

\begin{table}[H]
\caption{Alcance del proyecto}
\label{tab:alcance}
\begin{tabularx}{\linewidth}{Y Y}
\toprule
Incluye & No incluye \\
\midrule
Registro e inicio de sesión con correo o Google; JWT. & Recuperación de contraseña (el enlace de la pantalla de inicio no tiene endpoint). \\
Transacciones de ingreso y gasto, incluidas las recurrentes. & Inicio de sesión con Facebook (el botón existe, pero no está implementado). \\
Categorías predeterminadas y personalizadas. & Conexión con bancos o importación de extractos. \\
Presupuestos por categoría y período con alerta de superado. & Inversiones, créditos, contabilidad de empresas o impuestos. \\
Metas de ahorro con plan diario y progreso. & Varias monedas en una misma cuenta: cada usuario elige una. \\
Resumen, historial con filtros y gráficas. & Aplicación móvil nativa (la web es responsiva). \\
Perfil, moneda, foto, tema y umbral de alerta. & Notificaciones por correo o push. \\
Despliegue en la nube (Vercel, Render, Supabase). & Disponibilidad garantizada 24/7 (plan gratuito de Render; ver D-09). \\
\bottomrule
\end{tabularx}
\tablenote{Fuente: elaboración propia.}
\end{table}

El alcance es viable en el semestre: el producto se construyó entre abril y junio de 2026 y la fase de calidad entre agosto y noviembre de 2026, como muestra el cronograma de la sección \ref{sec:cronograma}.

\subsection{Valor frente a otras soluciones}

\begin{table}[H]
\caption{Comparación con la forma actual de controlar las finanzas}
\label{tab:valor}
\begin{tabularx}{\linewidth}{L{3.5cm} Y Y Y}
\toprule
Necesidad & Cuaderno & Hoja de cálculo & ShiftCash \\
\midrule
Registrar un gasto & A mano, sin validación & A mano; fórmulas propias & Formulario con validación en dos capas \\
Total por categoría & Suma manual & Fórmula que el usuario arma & Automático (resumen y gráficas) \\
Porcentaje del presupuesto & Cálculo manual & Fórmula y formato condicional propios & Automático, con estado «superado» y alerta \\
Gastos recurrentes & Se reescriben & Se copian & Se generan solos cada día \\
Plan de una meta & No existe & Debe diseñarse & Cuotas diarias automáticas y progreso \\
Acceso & Físico & Archivo local o en la nube & Web desde cualquier navegador, con sesión \\
\bottomrule
\end{tabularx}
\tablenote{Fuente: elaboración propia.}
\end{table}

Frente a la aplicación del banco, ShiftCash clasifica los movimientos según las categorías del propio usuario, reúne el efectivo y las distintas cuentas en un solo lugar y añade presupuestos y metas. Frente a otras aplicaciones de presupuesto, ShiftCash es gratuita para el usuario, está en español y redondea los montos a la moneda que él elige; \pendiente{comparación con dos o tres aplicaciones de presupuesto que el equipo haya probado (nombre, costo, idioma, funciones)}.

\paragraph{En cuánto mejora.} Para conocer el estado de un presupuesto con cuaderno u hoja de cálculo hay que sumar los gastos del período y dividirlos por el límite; en ShiftCash ese resultado aparece sin ninguna operación del usuario. El tiempo que toma registrar un gasto se mide con la tarea cronometrada PU-06 (Segunda parte): \pendiente{tiempo medio de registro con ShiftCash frente a cuaderno/hoja de cálculo, medido con 5 usuarios}.
